\documentclass[../main.tex]{subfiles}

\graphicspath{{../image/}{../figure/}}

\begin{document}

\chapter{平衡律方程组的 $BV$ 解}\label{chapter:ch16-平衡律组bv解}

此处的目的是讨论（可能非齐次的）严格双曲平衡律组的 Cauchy 问题 $BV$ 解的存在性与长时间行为。因而，本章可视为第 5.5 节的对应物，那里在经典解的语境下处理了相同的问题。
出于前几章给出的理由，研究将限制在单空间维的系统与小总变差初始数据上；然而，在这些限制之内，与第 5.5 节结果的类比相当地深远。于是，局部解的存在性将在对流与源项的温和限制下建立，而整体存在性将依赖于阻尼的存在。如同第 5.5 节那样，阻尼将由一个产生非负熵产生的耗散源诱导。

一个能在无任何辅助下产生整体解的强耗散源，在应用中很少遇到。在常见情形下，源仅仅部分耗散，整体解的存在性来自流与源之间的协同，编码于 Kawashima 条件中。

解可通过前几章给出的求解守恒律组的三种方法中的任一种构造，即随机选取、波前追踪或消失黏性，并结合算子分裂以计入源与/或非齐次性的效应。此处随机选取将是所选方法。

第 5.5 节中经典解的构造处于 $L^2$ 型 Sobolev 空间的范畴，由凸熵诱导的估计适合此目的。与之相反，此处的 $BV$ 理论需要 $L^1$ 型界。对于配有丰富熵族的两个平衡律组，$L^1$ 估计可通过构造在平衡态处具锥奇异的凸熵来导出。对于一般系统（其中所有可用熵函数在平衡态处光滑），$L^1$ 估计将由 $L^2$ 界导出，办法是利用双曲系统的有限依赖域性质。

\subfile{Chapter-16/section-01.tex}
\subfile{Chapter-16/section-02.tex}
\subfile{Chapter-16/section-03.tex}
\subfile{Chapter-16/section-04.tex}
\subfile{Chapter-16/section-05.tex}
\subfile{Chapter-16/section-06.tex}
\subfile{Chapter-16/section-07.tex}

\end{document}
